\chapter{Preliminaries and Notation}
\label{ch:preliminaries}

This chapter introduces the core concepts and notation used throughout this work.
The aim is to provide a common framework for interpreting the methods discussed
in the following chapters.

\section{Recurrent Neural Networks}

We consider recurrent neural networks (RNNs) that process sequences of inputs
$(x_t)_{t \geq 0}$ over discrete time steps. The network maintains a hidden
state $h_t \in \mathbb{R}^n$, updated according to
\begin{equation}
    h_t = f(h_{t-1}, x_t; \theta),
\end{equation}
where $\theta$ denotes the set of parameters. The network produces outputs
$y_t$ that depend on the hidden state.

\section{Spiking Neural Networks}

Spiking neural networks (SNNs) are a class of recurrent neural networks in
which neurons communicate through discrete binary events called \emph{spikes},
rather than continuous-valued activations. The output of neuron $j$ at time
$t$ is a binary variable $z^t_j \in \{0, 1\}$, where $z^t_j = 1$ indicates
that the neuron emits a spike at that timestep.

The most widely used model for spiking neurons is the leaky integrate-and-fire
(LIF) neuron. Each neuron maintains a scalar internal variable called the
\emph{membrane potential} $v^t_j$, which integrates incoming spikes from
connected neurons and decays over time. When $v^t_j$ exceeds a fixed threshold
$v_{th}$, the neuron emits a spike and the membrane potential is reset.

A key challenge arising from this formulation is that the spike generation
function is discontinuous, making standard gradient-based
optimization inapplicable. This issue and the techniques used to address it
are discussed in detail in Chapter~\ref{ch:sota}, in the context of the
learning algorithms that require it.

\section{Loss Functions and Temporal Dependencies}

We denote by $E$ a scalar loss function that depends on the sequence of
outputs. In recurrent systems, the effect of parameters at time $t$ may only
influence the loss at later times, leading to complex temporal dependencies
in the gradient $\frac{\mathrm{d}E}{\mathrm{d}\theta}$. Handling these 
dependencies efficiently is one of the central challenges in
training recurrent networks, and different methods address this problem in
different ways.

\section{Eligibility Traces}

Eligibility traces are variables associated with synaptic connections that
capture how past activity influences current updates. For a weight $W_{ji}$,
we denote the eligibility trace at time $t$ by $e^t_{ji}$. Intuitively, 
$e^t_{ji}$ measures the sensitivity of the current activity of
neuron $j$ to changes in $W_{ji}$, based only on locally available
information. 

In many methods, eligibility traces can be computed recursively
in forward time, making them well suited to online learning. Eligibility 
traces play a central role in online learning algorithms and provide an 
alternative mechanism for temporal credit assignment to the commonly used 
backpropagation variants.

\section{Reinforcement Learning}

In reinforcement learning (RL), an agent interacts with an environment over
discrete time steps with the goal of maximizing cumulative reward. At each
step $t$, the agent observes a state $S_t$, selects an action $A_t$, receives
a scalar reward $R_{t+1}$, and transitions to a new state $S_{t+1}$.

\paragraph{Policy.}
A \emph{policy} $\pi$ defines the agent's behavior as a mapping from states
to actions, either deterministic or stochastic. In the stochastic case,
$\pi(a \mid s)$ denotes the probability of selecting action $a$ in state $s$.
The objective of learning is to find a policy that maximizes expected return.

\paragraph{Return.}
The \emph{return} $G_t$ is the total discounted reward from time $t$ onward:
\begin{equation}
    G_t = \sum_{k=0}^{\infty} \gamma^k R_{t+1+k},
\end{equation}
where $\gamma \in [0,1)$ is a discount factor controlling the relative weight
of immediate versus future rewards.

\paragraph{Value function.}
Let $V_t \defeq V(S_t)$ denote the agent's estimate of the expected future
return from state $S_t$. Formally, the \emph{state-value function}
$v_\pi(s)$ measures the expected return starting from state $s$ under
policy $\pi$:
\begin{equation}
    v_\pi(s) = \mathbb{E}_\pi \left[ G_t \mid S_t = s \right].
\end{equation}
In practice, this is approximated by a learned estimate $V_t \approx
v_\pi(S_t)$, which the agent updates incrementally as it gains experience.

\paragraph{Bellman equations.}
The value function satisfies a recursive consistency condition known as the
\emph{Bellman equation}, which follows directly from expanding the definition
of $v_\pi(s)$ using the definition of $G_t$:
\begin{equation}
    v_\pi(s) = \mathbb{E}_\pi \left[ R_{t+1} + \gamma v_\pi(S_{t+1}) \mid S_t = s \right].
\end{equation}
This equation expresses the fundamental consistency requirement that the value of
a state must equal the expected immediate reward plus the discounted value
of the successor state. If our estimate $V_t$ were perfect, this condition
would hold exactly at every state. The degree to which it is violated at
any given timestep provides a natural signal for improving the estimate,
which is the basis of temporal-difference learning.

\paragraph{Temporal credit assignment.}
Because rewards may be delayed with respect to the actions that caused them,
the agent must determine which past decisions were responsible for the reward
currently being received. This is known as the \emph{temporal credit
assignment problem}, and is one of the central difficulties in reinforcement
learning, particularly in tasks with long time horizons.

\paragraph{Temporal-difference learning.}
Rather than waiting until the end of an episode to update value estimates,
temporal-difference (TD) learning exploits the Bellman equation directly.
If $V_t$ were a perfect estimate, the Bellman equation would require
$V_t = R_{t+1} + \gamma V_{t+1}$. The \emph{TD error} measures how far
the current estimate is from satisfying this condition:
\begin{equation}
    \delta_t = R_{t+1} + \gamma V_{t+1} - V_t,
\end{equation}
and provides an online, low-variance signal for updating both value estimates
and policy parameters. This quantity will appear repeatedly as a learning
signal throughout the following sections.

\clearpage
\section{Summary of Notation}

\begin{table}[h]
\centering
\begin{tabular}{ll}
\toprule
\textbf{Symbol} & \textbf{Meaning} \\
\midrule
$h_t$                       & Hidden state at time $t$ \\
$x_t$                       & Input at time $t$ \\
$z^t_j$                     & Spike output of neuron $j$ at time $t$ \\
$v^t_j$                     & Membrane potential of neuron $j$ at time $t$ \\
$v_{th}$                    & Spike threshold \\
$\alpha$                    & Membrane decay factor \\
$\mathcal{F}_\alpha$        & Exponential decay filter with parameter $\alpha$ \\
$W_{ji}$                    & Synaptic weight from neuron $i$ to $j$ \\
$e^t_{ji}$                  & Eligibility trace for synapse $ji$ at time $t$ \\
$L^t_j$                     & Learning signal for neuron $j$ at time $t$ \\
$E$                         & Loss function \\
$S_t,\ A_t,\ R_t$           & RL state, action, and reward at time $t$ \\
$\pi(a \mid s)$             & Stochastic policy \\
$G_t$                       & Discounted return at time $t$ \\
$\gamma$                    & Discount factor \\
$V_t$                       & Value estimate at time $t$ \\
$\delta_t$                  & Temporal-difference error \\
$\epsilon_j,\ \epsilon_{ji}$& Perturbation at neuron $j$ or synapse $ji$ \\
\bottomrule
\end{tabular}
\caption{Summary of notation used throughout this work.}
\end{table}
